\documentclass[border=10pt,tikz]{standalone}
\usepackage{tikz}
\usetikzlibrary{positioning,arrows.meta,shapes}
\usepackage{kotex}

% 출처: SEGGER, Unlimited Flash Breakpoints
% https://www.segger.com/products/debug-probes/j-link/technology/flash-breakpoints/
%  - hardware breakpoint를 먼저 쓰고, 다 쓰면 flash breakpoint로 전환
%  - flash sector를 다시 프로그래밍해 breakpoint를 설정·해제
%  - sector는 필요할 때만, 보통 대상 프로그램 실행을 시작할 때 프로그래밍
%  - 내장 명령어 집합 시뮬레이터와 프로그램 메모리 캐시로 flash 작업을 줄임

\begin{document}
% _design.tex — shared TikZ design system for all blog diagrams
% Inputs nothing on its own — meant to be \input{} from each diagram .tex
% AFTER \begin{document}, BEFORE \begin{tikzpicture}.
%
% Provides a unified visual language across GoF / DSA / EMC++ / Beautiful series:
%   - Colors: muted, accessible, color-coded by role
%   - Node styles: consistent sizing, rounded, subtle borders
%   - Edge styles: UML-style inheritance, aggregation, plus dashed alternates
%
% Usage:
%   \documentclass[border=10pt,tikz]{standalone}
%   \usepackage{tikz}
%   \usetikzlibrary{positioning,arrows.meta,shapes,fit,backgrounds}
%   \begin{document}
%   \input{../_design.tex}     % adjust relative path
%   \begin{tikzpicture}[blog]  % apply 'blog' style
%   ...

\definecolor{absfill}{RGB}{220,232,247}   % cool blue
\definecolor{absbord}{RGB}{72,118,182}
\definecolor{confill}{RGB}{249,222,222}   % warm red
\definecolor{conbord}{RGB}{197,93,93}
\definecolor{impfill}{RGB}{223,238,222}   % soft green
\definecolor{impbord}{RGB}{99,154,93}
\definecolor{clifill}{RGB}{251,239,205}   % pale yellow
\definecolor{clibord}{RGB}{200,167,72}
\definecolor{hifill} {RGB}{253,224,200}   % accent orange
\definecolor{hibord} {RGB}{210,135,68}
\definecolor{nodefill}{RGB}{250,250,250}
\definecolor{nodebord}{RGB}{120,120,120}
\definecolor{edgec}  {RGB}{80,80,80}
\definecolor{edgesc} {RGB}{145,145,145}
\definecolor{groupbord}{RGB}{200,200,200}

\tikzset{
  blog/.style={
    font=\sffamily\small,
    every node/.style={font=\sffamily\small}
  },
  % --- node roles ---
  cls/.style={
    rectangle, rounded corners=2pt,
    draw=nodebord, line width=0.4pt,
    minimum height=8mm, minimum width=28mm,
    inner sep=3pt, fill=nodefill, align=center
  },
  abs/.style={cls, draw=absbord, fill=absfill, font=\sffamily\itshape\small},
  con/.style={cls, draw=conbord, fill=confill},
  imp/.style={cls, draw=impbord, fill=impfill},
  cli/.style={cls, draw=clibord, fill=clifill},
  hi/.style={cls, draw=hibord, fill=hifill, font=\sffamily\bfseries\small},
  % --- circle (for trees / graph nodes) ---
  vx/.style={
    circle, draw=absbord, line width=0.4pt,
    minimum size=8mm, fill=absfill, inner sep=0pt
  },
  vxc/.style={vx, draw=conbord, fill=confill},
  vxh/.style={vx, draw=hibord, fill=hifill, font=\sffamily\bfseries},
  % --- decision diamond ---
  q/.style={
    diamond, draw=clibord, line width=0.4pt, fill=clifill,
    inner sep=2pt, aspect=2.4, align=center
  },
  % --- edges ---
  inh/.style={                                  % UML inheritance
    -{Triangle[length=3mm,width=3mm,open]},
    draw=edgec, line width=0.5pt
  },
  agg/.style={                                  % UML aggregation
    {Diamond[length=2.5mm,width=2.5mm,open]}-{Stealth[length=2.2mm]},
    draw=edgec, line width=0.45pt
  },
  uses/.style={                                 % directed reference
    -{Stealth[length=2.2mm]}, draw=edgec, line width=0.5pt
  },
  arr/.style={                                  % plain arrow (graphs/trees)
    -{Stealth[length=2mm]}, draw=edgec, line width=0.45pt
  },
  edge/.style={                                 % undirected (trees)
    draw=edgec, line width=0.45pt
  },
  alt/.style={                                  % dashed: similar/alt
    -{Stealth[length=2mm]}, dashed, draw=edgesc, line width=0.45pt
  },
  ret/.style={                                  % dashed creates/returns
    -{Stealth[length=2mm]}, dashed, draw=edgesc, line width=0.45pt
  },
  thr/.style={                                  % threading link (red dashed)
    ->, dashed, draw=conbord, line width=0.5pt
  },
  % --- group / region ---
  group/.style={
    rectangle, rounded corners=4pt,
    draw=groupbord, line width=0.4pt,
    inner sep=10pt, fill=black!2
  },
  glabel/.style={font=\sffamily\bfseries\small, anchor=north west, inner sep=2pt}
}

% Korean font convenience macro — included by every Hangul diagram
\providecommand{\KOR}{}


\begin{tikzpicture}[blog,
    box/.style={rectangle, rounded corners=3pt, draw=absbord, fill=absfill,
                text width=52mm, minimum height=12mm, align=center,
                font=\sffamily\small},
    hwbox/.style={box, draw=impbord, fill=impfill},
    flbox/.style={box, draw=hibord, fill=hifill},
    dec/.style={diamond, aspect=2.4, draw=nodebord, fill=nodefill, align=center,
                font=\sffamily\small, inner sep=1pt},
    note/.style={rectangle, rounded corners=2pt, draw=groupbord, fill=white,
                 align=left, font=\sffamily\footnotesize, inner sep=5pt},
    arr/.style={-{Stealth[length=2.4mm]}, draw=edgec, line width=0.55pt},
    lbl/.style={font=\sffamily\footnotesize, text=inkc, fill=white, inner sep=2pt}
]

\node[box] (req) at (0, 4.2)
  {GDB가 breakpoint 설정 요청\\[2pt]\footnotesize J-Link 소프트웨어가 처리};

\node[dec] (q) at (0, 1.9) {hardware breakpoint\\[2pt]남아 있는가?};

\node[hwbox] (hw) at (-4.2, -0.6)
  {\textbf{hardware breakpoint}\\[2pt]\footnotesize 먼저 사용};
\node[flbox] (fl) at (4.2, -0.6)
  {\textbf{flash breakpoint}\\[2pt]\footnotesize flash sector를 다시 프로그래밍해\\[2pt]\footnotesize breakpoint를 설정·해제};

\draw[arr] (req) -- (q);
\draw[arr] (q.west) -| node[lbl, pos=0.25, above]{예} (hw.north);
\draw[arr] (q.east) -| node[lbl, pos=0.5, above]{아니오 (모두 사용 중)} (fl.north);

\node[note, below=6mm of fl] (n1)
  {$\bullet$ sector는 필요할 때만 프로그래밍\\[2pt]\hphantom{$\bullet$ }(보통 대상 프로그램 실행을 시작할 때)\\[2pt]
   $\bullet$ 내장 명령어 집합 시뮬레이터와\\[2pt]\hphantom{$\bullet$ }프로그램 메모리 캐시로 flash 작업을 줄임};
\draw[draw=groupbord, line width=0.5pt] (fl.south) -- (n1.north);

\node[font=\sffamily\footnotesize, text=mutec, align=center] at (0, -4.6)
  {SEGGER 설명 기준. 지원 장치·flash 영역과 라이선스(J-Link 단위)를 확인합니다.};

\end{tikzpicture}
\end{document}
